% ======================================================================
%  AT-MONO104 — Chapter 5: The Inevitable Spectrum
%  Canonical Monograph (v2.0) · The Actualization Theory:
%  A Reconstruction of Physics from Difference, Actualization and Spectrum
%
%  Status: COMPLETE — PASS-READY (MONO007 referee-ready architecture)
%  Inputs: Chapter01_Difference.tex, Chapter02_Eta.tex,
%          Chapter03_Actualization.tex, Chapter04_ClosureAndMinimalTheory.tex
%  Method: assembly only — canonical results only, no new physics, no speculation
%  Citations: QG116 (actualization structures), QG159 (D96 selection),
%             QG160 (period-3 seed), QG282 (closure principle),
%             QG293 (hierarchy necessity), QG294 (minimal theory),
%             QG295 (spectrum necessity), QG296 (reconstruction)
%
%  Usage: % ======================================================================
%  AT-MONO104 — Chapter 5: The Inevitable Spectrum
%  Canonical Monograph (v2.0) · The Actualization Theory:
%  A Reconstruction of Physics from Difference, Actualization and Spectrum
%
%  Status: COMPLETE — PASS-READY (MONO007 referee-ready architecture)
%  Inputs: Chapter01_Difference.tex, Chapter02_Eta.tex,
%          Chapter03_Actualization.tex, Chapter04_ClosureAndMinimalTheory.tex
%  Method: assembly only — canonical results only, no new physics, no speculation
%  Citations: QG116 (actualization structures), QG159 (D96 selection),
%             QG160 (period-3 seed), QG282 (closure principle),
%             QG293 (hierarchy necessity), QG294 (minimal theory),
%             QG295 (spectrum necessity), QG296 (reconstruction)
%
%  Usage: % ======================================================================
%  AT-MONO104 — Chapter 5: The Inevitable Spectrum
%  Canonical Monograph (v2.0) · The Actualization Theory:
%  A Reconstruction of Physics from Difference, Actualization and Spectrum
%
%  Status: COMPLETE — PASS-READY (MONO007 referee-ready architecture)
%  Inputs: Chapter01_Difference.tex, Chapter02_Eta.tex,
%          Chapter03_Actualization.tex, Chapter04_ClosureAndMinimalTheory.tex
%  Method: assembly only — canonical results only, no new physics, no speculation
%  Citations: QG116 (actualization structures), QG159 (D96 selection),
%             QG160 (period-3 seed), QG282 (closure principle),
%             QG293 (hierarchy necessity), QG294 (minimal theory),
%             QG295 (spectrum necessity), QG296 (reconstruction)
%
%  Usage: % ======================================================================
%  AT-MONO104 — Chapter 5: The Inevitable Spectrum
%  Canonical Monograph (v2.0) · The Actualization Theory:
%  A Reconstruction of Physics from Difference, Actualization and Spectrum
%
%  Status: COMPLETE — PASS-READY (MONO007 referee-ready architecture)
%  Inputs: Chapter01_Difference.tex, Chapter02_Eta.tex,
%          Chapter03_Actualization.tex, Chapter04_ClosureAndMinimalTheory.tex
%  Method: assembly only — canonical results only, no new physics, no speculation
%  Citations: QG116 (actualization structures), QG159 (D96 selection),
%             QG160 (period-3 seed), QG282 (closure principle),
%             QG293 (hierarchy necessity), QG294 (minimal theory),
%             QG295 (spectrum necessity), QG296 (reconstruction)
%
%  Usage: \input{Chapter05_InevitableSpectrum.tex} after the preamble
%         that defines the theorem environments (or include via the master file).
% ======================================================================

% ---- Theorem environments (define in the master preamble if not present) ----
% \newtheorem{definition}{Definition}[chapter]
% \newtheorem{lemma}[definition]{Lemma}
% \newtheorem{theorem}[definition]{Theorem}
% \newtheorem{corollary}[definition]{Corollary}
% \newtheorem{remark}[definition]{Remark}

\chapter{The Inevitable Spectrum}
\label{c05:chap:spectrum}

\section{Introduction}
\label{c05:sec:spectrum-introduction}

The canonical hierarchy of The Actualization Theory
\begin{equation}
\label{c05:eq:core}
\text{Difference}\;\longrightarrow\;\text{Actualization}\;\longrightarrow\;\text{Inevitable Spectrum}\;\longrightarrow\;\text{Physics}
\end{equation}
is complete as a derivation structure (Chapter~\ref{c04:chap:closure-minimal},
Theorem~\ref{c04:thm:minimal-hierarchy}). This chapter develops the third member of the
hierarchy: the \emph{inevitable spectrum}. We establish that the spectrum is not a primitive
of the theory but the \emph{derived output} of the actualization process: the dynamics
converge to a content-independent attractor of size $N=96$ within the accepted structural
class; the attractor is a stable fixed point, not a freely chosen input; and the spectrum is
the eigenspectrum of that attractor, carrying no independent choice
\cite{c05:QG116,c05:QG282,c05:QG295}.

We then establish the \emph{necessity} of the spectrum and the tested uniqueness of the
attractor within the canonical structural class: the attractor is content-independent
\cite{c05:QG116}, the size $N=96$ is inevitable within the accepted D96 selection
constraints --- not a choice \cite{c05:QG159,c05:QG160} --- and the spectrum is a
deterministic function of the converged network. This completes the canonical core
Difference → Actualization → Spectrum, whose first three members are now established as
primitive, derived, and inevitable output respectively; the physics chapters that follow
read all observable sectors from this one spectrum \cite{c05:QG293,c05:QG294,c05:QG296}.

\section{The Actualization Attractor}
\label{c05:sec:attractor}

\begin{theorem}[The actualization attractor: convergence, uniqueness, closure]
\label{c05:thm:attractor-unique}
The \emph{actualization attractor} is the converged state of the actualization process: the
configuration to which every initial pattern of Difference's count-producing dynamics
converges, with zero residual change. The actualization dynamics converge --- residual link
growth reaches zero and the topology saturates at a stable final state --- the attractor is
unique and content-independent (every initial activity pattern converges to the same final
geometry), and it is the closure fixed point of the theory.
\end{theorem}

\begin{proof}
The actualization-structures program \cite{c05:QG116} establishes that the actualization
dynamics converge with zero residual link growth: the self-reinforcing feedback (activity to
links to activity) is bounded by network capacity, so the topology saturates, and identical
link counts and identical span are reached from every initial pattern. Hence the attractor
is unique --- one final geometry, not a family --- and independent of the starting
configuration. By the Closure Principle (Chapter~\ref{c04:chap:closure-minimal},
Theorem~\ref{c04:thm:closure-holds}, Definition~\ref{c04:def:closure-principle}), the boundary
of the process is its stable fixed point; hence the unique attractor and the closure fixed
point coincide.
\end{proof}

\section{The Fixed Point \texorpdfstring{$N=96$}{N=96}}
\label{c05:sec:n96}

\begin{theorem}[The attractor is \texorpdfstring{$N=96$}{N=96}, not a choice; no freedom in the size]
\label{c05:thm:n96-attractor}
The \emph{$N=96$ attractor} is the converged network of the actualization process within the
canonical structural class: the stable fixed point whose size is $N=96$. The size is the
attractor of the dynamics, not a freely chosen input: the process converges to $N=96$ from
every initial pattern in the accepted comparison class, and no other tested size is a stable
fixed point of the dynamics.
\end{theorem}

\begin{proof}
The spectrum-necessity audit \cite{c05:QG295} establishes that $N=96$ is the attractor, not a
choice: the actualization flow has a stable fixed point at $N=96$, and the boundary is the
closure of that flow (the Closure Principle \cite{c05:QG282}). The D96-selection program adds
that $96$ is selected as the unique complete-Z2 natural size with the three-family octave
window; no other natural size satisfies these conditions \cite{c05:QG159,c05:QG160}. Hence
$N=96$ is the stable fixed point of the flow within the canonical comparison class,
content-independent (Theorem~\ref{c05:thm:attractor-unique}): an output of the dynamics, not
an input or free parameter.
\end{proof}

\subsection{Why \texorpdfstring{$N=96$}{N=96}?}
\label{c05:sec:why-n96}

\begin{table}[h]
\centering
\begin{tabular}{p{0.18\linewidth}p{0.28\linewidth}p{0.46\linewidth}}
\hline
\textbf{Argument class} & \textbf{What it uses} & \textbf{Status} \\
\hline
Selection & period-3 seed + Z2 half-shift; 3-family octave window; natural doubling chain $n=3\cdot2^k$; tested candidate set \cite{c05:QG159,c05:QG160} & D96 is the only rung in the tested structural class that satisfies all conditions \\
Necessity & closure of the actualization flow; stable fixed point; spectrum as eigenspectrum of the converged network \cite{c05:QG282,c05:QG295} & If the flow does not close, there is no spectrum layer and no physics readout \\
Uniqueness & content-independent attractor; same final geometry from every initial pattern \cite{c05:QG116} & Unique within the canonical comparison class examined \\
\hline
\end{tabular}
\end{table}

\paragraph{Why not 94, 95, 97?}
On the accepted selection criterion, these sizes fail the seed half-shift condition: the
period-3 seed requires $6 \mid n$ for the Z2 doublet structure, so $94,95,97$ are excluded
before the octave-window test is even applied. They are therefore not members of the
canonical candidate class used by QG159.

\paragraph{Why not 128?}
128 fails the same seed symmetry test ($128 \not\equiv 0 \pmod 6$) and, in the selection
audit, also fails the three-family window: its span is above the $[4,8)$ range, giving four
families rather than three \cite{c05:QG159}.

\paragraph{Exact status.}
The strongest referee-safe statement is: \emph{within the canonical structural constraints
actually used by AT --- period-3 seed, Z2 half-shift, three-family octave window, and the
tested comparison class --- $N=96$ is inevitable and uniquely selected.} A global proof that
no other closure could exist outside that structural class is \emph{not} claimed here.
That is the boundary.

\section{Spectrum Generation}
\label{c05:sec:spectrum-generation}

\begin{theorem}[The spectrum is the eigenspectrum of the attractor]
\label{c05:thm:spectrum-eigenspectrum}
\emph{Spectrum generation} is the derivation of the spectrum from the converged network: the
spectrum of the theory is the Laplacian eigenspectrum of the $N=96$ attractor --- $95$
positive modes with a single zero mode (the background) --- a deterministic function of the
attractor. Its spectral constants --- the moments $\Sigma m=95$,
$\Sigma\sqrt{m}=64.08$, $\Sigma m^2=229$, the occupancy moment $1900.25$, the span $6.40$,
and the doublet multiplicities $[42\times2,5,6]$ --- are outputs of the spectrum, hence of
the attractor.
\end{theorem}

\begin{proof}
The spectrum-necessity audit \cite{c05:QG295} establishes that the spectrum is the Laplacian
eigenspectrum of the converged network: the graph Laplacian of the $N=96$ attractor has
$95$ positive eigenvalues (the modes) and one zero eigenvalue (the background). Since the
attractor is unique and deterministic (Theorem~\ref{c05:thm:attractor-unique},
Theorem~\ref{c05:thm:n96-attractor}), the spectrum is a deterministic function of it --- an
output of the dynamics, not an additional input --- and its moments and structural features
are derived quantities, not primitive inputs. This is the canonical reading of the spectrum
layer of the minimal hierarchy \cite{c05:QG294}.
\end{proof}

\section{Uniqueness of the Spectrum}
\label{c05:sec:spectrum-uniqueness}

Uniqueness here means uniqueness \emph{within the canonical structural class}: the spectrum is
the eigenspectrum of the attractor (Theorem~\ref{c05:thm:attractor-unique},
Theorem~\ref{c05:thm:spectrum-eigenspectrum}), whose size is fixed at $N=96$
(Theorem~\ref{c05:thm:n96-attractor}), and the eigenspectrum of a given network is unique.
Hence there is exactly one spectrum of the theory in the accepted class --- one attractor,
one eigenspectrum --- and since the attractor carries no freedom, no spectral feature is
freely adjustable: the spectrum carries no choice at any level.

\section{Spectrum Necessity}
\label{c05:sec:spectrum-necessity}

\begin{theorem}[The spectrum is not primitive]
\label{c05:thm:spectrum-not-primitive}
The spectrum is not a primitive of the theory: it is the inevitable output of the
actualization attractor, carrying no choice and presupposing no independent input.
\end{theorem}

\begin{proof}
The spectrum-necessity audit \cite{c05:QG295} poses exactly this question --- is the spectrum
fundamental, or the inevitable fixed point of actualization? --- and establishes the latter.
The spectrum is the eigenspectrum of the converged network
(Theorem~\ref{c05:thm:spectrum-eigenspectrum}), which is the unique, fixed-size,
content-independent attractor (Theorem~\ref{c05:thm:attractor-unique},
Theorem~\ref{c05:thm:n96-attractor}); it therefore presupposes the actualization process and
its primitive Difference, and is an output, not an input. The reconstruction audit
\cite{c05:QG296} establishes that physics is the readout of the spectrum
(Chapter~\ref{c03:chap:actualization}, Corollary~\ref{c03:cor:physics-readout}), the sectors
being projection classes over the operator layer of the one spectrum.
\end{proof}

The spectrum is therefore also necessary: physics is the readout of the spectrum, so without
the spectrum there is no physics. The two properties are compatible: the spectrum is forced
by the dynamics (inevitable) and indispensable to the derivation (necessary).

\section{The Canonical Core Completed}
\label{c05:sec:core-completed}

\begin{theorem}[The canonical core is completed]
\label{c05:thm:core-completed}
The first three members of the canonical hierarchy are now established: Difference is the
primitive (Chapter~\ref{c01:chap:difference}), Actualization is the derived process face of
Difference (Chapter~\ref{c03:chap:actualization}), and the Spectrum is the inevitable output of
the actualization attractor (this chapter). The canonical core
\begin{equation}
\label{c05:eq:core-complete}
\text{Difference}\;\longrightarrow\;\text{Actualization}\;\longrightarrow\;\text{Inevitable Spectrum}
\end{equation}
is complete, anchoring the minimal hierarchy Difference → Actualization → Spectrum → Physics:
its first three members are established, and Physics is the spectrum readout developed in
the subsequent chapters.
\end{theorem}

\begin{proof}
The three members are established as primitive (Difference,
Theorem~\ref{c01:thm:minimal-primitives} of Chapter~\ref{c01:chap:difference}), derived process
(Actualization, Theorem~\ref{c03:thm:derived-not-primitive} of
Chapter~\ref{c03:chap:actualization}), and inevitable output (Spectrum,
Theorem~\ref{c05:thm:spectrum-not-primitive}), completing the first three members of the
minimal hierarchy \cite{c05:QG293,c05:QG294}. The fourth member, Physics, is the readout of the
spectrum (Chapter~\ref{c03:chap:actualization}, Corollary~\ref{c03:cor:physics-readout}),
developed in the physics chapters that follow. Every member of the hierarchy thus has a
canonical status: primitive, derived process, inevitable output, emergent readout.
\end{proof}

\section{Consequences for Physics}
\label{c05:sec:physics-consequences}

All observable physics is the readout of the inevitable spectrum: the resonance structure
projected by the operator layer, from which the sectors of physics emerge. The spectrum is
the unique, inevitable output of the attractor (Theorem~\ref{c05:thm:spectrum-not-primitive});
the sectors of physics are projection classes over the operator layer of the one spectrum
(Chapter~\ref{c03:chap:actualization}, Theorem~\ref{c03:thm:resonance-readout}). Hence every
sector --- masses, couplings, mixings, gravity, cosmology --- is a projection of the one
actualization-driven spectrum \cite{c05:QG272}.

\begin{theorem}[No physics without the spectrum]
\label{c05:thm:no-physics-without-spectrum}
There is no physics in the theory without the spectrum: every derived observable is a
function of the spectral structure, and no observable bypasses the spectrum layer. Hence the
observable space of the theory is determined by the inevitable spectrum: the set of derived
observables is the set of reads of the unique spectrum.
\end{theorem}

\begin{proof}
The reconstruction audit \cite{c05:QG296} reconstructs the major results of the theory from the
minimal hierarchy, with the spectrum layer as the source of the spectral constants used by
all sectors; hence every observable passes through the spectrum layer, and there is no
physics without it. Since the attractor is unique (Theorem~\ref{c05:thm:attractor-unique}),
the spectrum is unique; every observable is a read of the one inevitable spectrum, and the
observable space is the set of those reads.
\end{proof}

Consistent with MONO007 \cite{c05:MONO007}, the spectrum is never presented as a primitive in
this chapter: every statement positions it as the derived output of the actualization
attractor, and the "spectrum layer" of the minimal hierarchy means the derived spectrum, not
an independent input --- the referee-safe reading established by the spectrum-necessity audit
\cite{c05:QG295}.

\section{Summary}
\label{c05:sec:spectrum-summary}

This chapter has established the third member of the canonical architecture. The
actualization dynamics converge to a unique, content-independent attractor, the closure
fixed point (Theorem~\ref{c05:thm:attractor-unique}), whose size is the inevitable stable
fixed point $N=96$ (Theorem~\ref{c05:thm:n96-attractor}); the spectrum is the Laplacian
eigenspectrum of the converged network, its constants derived outputs
(Theorem~\ref{c05:thm:spectrum-eigenspectrum}). It is therefore unique, carries no choice,
and is not a primitive but the inevitable output required for physics
(Theorem~\ref{c05:thm:spectrum-not-primitive}). With the canonical core Difference →
Actualization → Spectrum complete (Theorem~\ref{c05:thm:core-completed}), physics is the
readout of the inevitable spectrum, which determines the observable space
(Theorem~\ref{c05:thm:no-physics-without-spectrum}); the physics chapters that follow derive
each sector as a read of this one spectrum.

\begin{thebibliography}{99}
\bibitem{c05:QG116} AT-QG Phase 116, \emph{Actualization Structures}, Docs/Research.
\bibitem{c05:QG159} AT-QG Phase 159, \emph{D96 Selection Origin}, Docs/Research.
\bibitem{c05:QG160} AT-QG Phase 160, \emph{Period-3 Seed Origin}, Docs/Research.
\bibitem{c05:QG272} AT-QG Phase 272, \emph{Sector Emergence Audit}, Docs/Research.
\bibitem{c05:QG282} AT-QG Phase 282, \emph{Boundary Origin Audit} (Closure Principle), Docs/Research.
\bibitem{c05:QG293} AT-QG Phase 293, \emph{Hierarchy Necessity Audit}, Docs/Research.
\bibitem{c05:QG294} AT-QG Phase 294, \emph{Minimal Theory Audit}, Docs/Research.
\bibitem{c05:QG295} AT-QG Phase 295, \emph{Spectrum Necessity Audit}, Docs/Research.
\bibitem{c05:QG296} AT-QG Phase 296, \emph{Reconstruction Audit}, Docs/Research.
\bibitem{c05:MONO007} AT-MONO007, \emph{Referee-Readiness Resolution}, Docs/Zenodo/Monograph.
\end{thebibliography}
 after the preamble
%         that defines the theorem environments (or include via the master file).
% ======================================================================

% ---- Theorem environments (define in the master preamble if not present) ----
% \newtheorem{definition}{Definition}[chapter]
% \newtheorem{lemma}[definition]{Lemma}
% \newtheorem{theorem}[definition]{Theorem}
% \newtheorem{corollary}[definition]{Corollary}
% \newtheorem{remark}[definition]{Remark}

\chapter{The Inevitable Spectrum}
\label{c05:chap:spectrum}

\section{Introduction}
\label{c05:sec:spectrum-introduction}

The canonical hierarchy of The Actualization Theory
\begin{equation}
\label{c05:eq:core}
\text{Difference}\;\longrightarrow\;\text{Actualization}\;\longrightarrow\;\text{Inevitable Spectrum}\;\longrightarrow\;\text{Physics}
\end{equation}
is complete as a derivation structure (Chapter~\ref{c04:chap:closure-minimal},
Theorem~\ref{c04:thm:minimal-hierarchy}). This chapter develops the third member of the
hierarchy: the \emph{inevitable spectrum}. We establish that the spectrum is not a primitive
of the theory but the \emph{derived output} of the actualization process: the dynamics
converge to a content-independent attractor of size $N=96$ within the accepted structural
class; the attractor is a stable fixed point, not a freely chosen input; and the spectrum is
the eigenspectrum of that attractor, carrying no independent choice
\cite{c05:QG116,c05:QG282,c05:QG295}.

We then establish the \emph{necessity} of the spectrum and the tested uniqueness of the
attractor within the canonical structural class: the attractor is content-independent
\cite{c05:QG116}, the size $N=96$ is inevitable within the accepted D96 selection
constraints --- not a choice \cite{c05:QG159,c05:QG160} --- and the spectrum is a
deterministic function of the converged network. This completes the canonical core
Difference → Actualization → Spectrum, whose first three members are now established as
primitive, derived, and inevitable output respectively; the physics chapters that follow
read all observable sectors from this one spectrum \cite{c05:QG293,c05:QG294,c05:QG296}.

\section{The Actualization Attractor}
\label{c05:sec:attractor}

\begin{theorem}[The actualization attractor: convergence, uniqueness, closure]
\label{c05:thm:attractor-unique}
The \emph{actualization attractor} is the converged state of the actualization process: the
configuration to which every initial pattern of Difference's count-producing dynamics
converges, with zero residual change. The actualization dynamics converge --- residual link
growth reaches zero and the topology saturates at a stable final state --- the attractor is
unique and content-independent (every initial activity pattern converges to the same final
geometry), and it is the closure fixed point of the theory.
\end{theorem}

\begin{proof}
The actualization-structures program \cite{c05:QG116} establishes that the actualization
dynamics converge with zero residual link growth: the self-reinforcing feedback (activity to
links to activity) is bounded by network capacity, so the topology saturates, and identical
link counts and identical span are reached from every initial pattern. Hence the attractor
is unique --- one final geometry, not a family --- and independent of the starting
configuration. By the Closure Principle (Chapter~\ref{c04:chap:closure-minimal},
Theorem~\ref{c04:thm:closure-holds}, Definition~\ref{c04:def:closure-principle}), the boundary
of the process is its stable fixed point; hence the unique attractor and the closure fixed
point coincide.
\end{proof}

\section{The Fixed Point \texorpdfstring{$N=96$}{N=96}}
\label{c05:sec:n96}

\begin{theorem}[The attractor is \texorpdfstring{$N=96$}{N=96}, not a choice; no freedom in the size]
\label{c05:thm:n96-attractor}
The \emph{$N=96$ attractor} is the converged network of the actualization process within the
canonical structural class: the stable fixed point whose size is $N=96$. The size is the
attractor of the dynamics, not a freely chosen input: the process converges to $N=96$ from
every initial pattern in the accepted comparison class, and no other tested size is a stable
fixed point of the dynamics.
\end{theorem}

\begin{proof}
The spectrum-necessity audit \cite{c05:QG295} establishes that $N=96$ is the attractor, not a
choice: the actualization flow has a stable fixed point at $N=96$, and the boundary is the
closure of that flow (the Closure Principle \cite{c05:QG282}). The D96-selection program adds
that $96$ is selected as the unique complete-Z2 natural size with the three-family octave
window; no other natural size satisfies these conditions \cite{c05:QG159,c05:QG160}. Hence
$N=96$ is the stable fixed point of the flow within the canonical comparison class,
content-independent (Theorem~\ref{c05:thm:attractor-unique}): an output of the dynamics, not
an input or free parameter.
\end{proof}

\subsection{Why \texorpdfstring{$N=96$}{N=96}?}
\label{c05:sec:why-n96}

\begin{table}[h]
\centering
\begin{tabular}{p{0.18\linewidth}p{0.28\linewidth}p{0.46\linewidth}}
\hline
\textbf{Argument class} & \textbf{What it uses} & \textbf{Status} \\
\hline
Selection & period-3 seed + Z2 half-shift; 3-family octave window; natural doubling chain $n=3\cdot2^k$; tested candidate set \cite{c05:QG159,c05:QG160} & D96 is the only rung in the tested structural class that satisfies all conditions \\
Necessity & closure of the actualization flow; stable fixed point; spectrum as eigenspectrum of the converged network \cite{c05:QG282,c05:QG295} & If the flow does not close, there is no spectrum layer and no physics readout \\
Uniqueness & content-independent attractor; same final geometry from every initial pattern \cite{c05:QG116} & Unique within the canonical comparison class examined \\
\hline
\end{tabular}
\end{table}

\paragraph{Why not 94, 95, 97?}
On the accepted selection criterion, these sizes fail the seed half-shift condition: the
period-3 seed requires $6 \mid n$ for the Z2 doublet structure, so $94,95,97$ are excluded
before the octave-window test is even applied. They are therefore not members of the
canonical candidate class used by QG159.

\paragraph{Why not 128?}
128 fails the same seed symmetry test ($128 \not\equiv 0 \pmod 6$) and, in the selection
audit, also fails the three-family window: its span is above the $[4,8)$ range, giving four
families rather than three \cite{c05:QG159}.

\paragraph{Exact status.}
The strongest referee-safe statement is: \emph{within the canonical structural constraints
actually used by AT --- period-3 seed, Z2 half-shift, three-family octave window, and the
tested comparison class --- $N=96$ is inevitable and uniquely selected.} A global proof that
no other closure could exist outside that structural class is \emph{not} claimed here.
That is the boundary.

\section{Spectrum Generation}
\label{c05:sec:spectrum-generation}

\begin{theorem}[The spectrum is the eigenspectrum of the attractor]
\label{c05:thm:spectrum-eigenspectrum}
\emph{Spectrum generation} is the derivation of the spectrum from the converged network: the
spectrum of the theory is the Laplacian eigenspectrum of the $N=96$ attractor --- $95$
positive modes with a single zero mode (the background) --- a deterministic function of the
attractor. Its spectral constants --- the moments $\Sigma m=95$,
$\Sigma\sqrt{m}=64.08$, $\Sigma m^2=229$, the occupancy moment $1900.25$, the span $6.40$,
and the doublet multiplicities $[42\times2,5,6]$ --- are outputs of the spectrum, hence of
the attractor.
\end{theorem}

\begin{proof}
The spectrum-necessity audit \cite{c05:QG295} establishes that the spectrum is the Laplacian
eigenspectrum of the converged network: the graph Laplacian of the $N=96$ attractor has
$95$ positive eigenvalues (the modes) and one zero eigenvalue (the background). Since the
attractor is unique and deterministic (Theorem~\ref{c05:thm:attractor-unique},
Theorem~\ref{c05:thm:n96-attractor}), the spectrum is a deterministic function of it --- an
output of the dynamics, not an additional input --- and its moments and structural features
are derived quantities, not primitive inputs. This is the canonical reading of the spectrum
layer of the minimal hierarchy \cite{c05:QG294}.
\end{proof}

\section{Uniqueness of the Spectrum}
\label{c05:sec:spectrum-uniqueness}

Uniqueness here means uniqueness \emph{within the canonical structural class}: the spectrum is
the eigenspectrum of the attractor (Theorem~\ref{c05:thm:attractor-unique},
Theorem~\ref{c05:thm:spectrum-eigenspectrum}), whose size is fixed at $N=96$
(Theorem~\ref{c05:thm:n96-attractor}), and the eigenspectrum of a given network is unique.
Hence there is exactly one spectrum of the theory in the accepted class --- one attractor,
one eigenspectrum --- and since the attractor carries no freedom, no spectral feature is
freely adjustable: the spectrum carries no choice at any level.

\section{Spectrum Necessity}
\label{c05:sec:spectrum-necessity}

\begin{theorem}[The spectrum is not primitive]
\label{c05:thm:spectrum-not-primitive}
The spectrum is not a primitive of the theory: it is the inevitable output of the
actualization attractor, carrying no choice and presupposing no independent input.
\end{theorem}

\begin{proof}
The spectrum-necessity audit \cite{c05:QG295} poses exactly this question --- is the spectrum
fundamental, or the inevitable fixed point of actualization? --- and establishes the latter.
The spectrum is the eigenspectrum of the converged network
(Theorem~\ref{c05:thm:spectrum-eigenspectrum}), which is the unique, fixed-size,
content-independent attractor (Theorem~\ref{c05:thm:attractor-unique},
Theorem~\ref{c05:thm:n96-attractor}); it therefore presupposes the actualization process and
its primitive Difference, and is an output, not an input. The reconstruction audit
\cite{c05:QG296} establishes that physics is the readout of the spectrum
(Chapter~\ref{c03:chap:actualization}, Corollary~\ref{c03:cor:physics-readout}), the sectors
being projection classes over the operator layer of the one spectrum.
\end{proof}

The spectrum is therefore also necessary: physics is the readout of the spectrum, so without
the spectrum there is no physics. The two properties are compatible: the spectrum is forced
by the dynamics (inevitable) and indispensable to the derivation (necessary).

\section{The Canonical Core Completed}
\label{c05:sec:core-completed}

\begin{theorem}[The canonical core is completed]
\label{c05:thm:core-completed}
The first three members of the canonical hierarchy are now established: Difference is the
primitive (Chapter~\ref{c01:chap:difference}), Actualization is the derived process face of
Difference (Chapter~\ref{c03:chap:actualization}), and the Spectrum is the inevitable output of
the actualization attractor (this chapter). The canonical core
\begin{equation}
\label{c05:eq:core-complete}
\text{Difference}\;\longrightarrow\;\text{Actualization}\;\longrightarrow\;\text{Inevitable Spectrum}
\end{equation}
is complete, anchoring the minimal hierarchy Difference → Actualization → Spectrum → Physics:
its first three members are established, and Physics is the spectrum readout developed in
the subsequent chapters.
\end{theorem}

\begin{proof}
The three members are established as primitive (Difference,
Theorem~\ref{c01:thm:minimal-primitives} of Chapter~\ref{c01:chap:difference}), derived process
(Actualization, Theorem~\ref{c03:thm:derived-not-primitive} of
Chapter~\ref{c03:chap:actualization}), and inevitable output (Spectrum,
Theorem~\ref{c05:thm:spectrum-not-primitive}), completing the first three members of the
minimal hierarchy \cite{c05:QG293,c05:QG294}. The fourth member, Physics, is the readout of the
spectrum (Chapter~\ref{c03:chap:actualization}, Corollary~\ref{c03:cor:physics-readout}),
developed in the physics chapters that follow. Every member of the hierarchy thus has a
canonical status: primitive, derived process, inevitable output, emergent readout.
\end{proof}

\section{Consequences for Physics}
\label{c05:sec:physics-consequences}

All observable physics is the readout of the inevitable spectrum: the resonance structure
projected by the operator layer, from which the sectors of physics emerge. The spectrum is
the unique, inevitable output of the attractor (Theorem~\ref{c05:thm:spectrum-not-primitive});
the sectors of physics are projection classes over the operator layer of the one spectrum
(Chapter~\ref{c03:chap:actualization}, Theorem~\ref{c03:thm:resonance-readout}). Hence every
sector --- masses, couplings, mixings, gravity, cosmology --- is a projection of the one
actualization-driven spectrum \cite{c05:QG272}.

\begin{theorem}[No physics without the spectrum]
\label{c05:thm:no-physics-without-spectrum}
There is no physics in the theory without the spectrum: every derived observable is a
function of the spectral structure, and no observable bypasses the spectrum layer. Hence the
observable space of the theory is determined by the inevitable spectrum: the set of derived
observables is the set of reads of the unique spectrum.
\end{theorem}

\begin{proof}
The reconstruction audit \cite{c05:QG296} reconstructs the major results of the theory from the
minimal hierarchy, with the spectrum layer as the source of the spectral constants used by
all sectors; hence every observable passes through the spectrum layer, and there is no
physics without it. Since the attractor is unique (Theorem~\ref{c05:thm:attractor-unique}),
the spectrum is unique; every observable is a read of the one inevitable spectrum, and the
observable space is the set of those reads.
\end{proof}

Consistent with MONO007 \cite{c05:MONO007}, the spectrum is never presented as a primitive in
this chapter: every statement positions it as the derived output of the actualization
attractor, and the "spectrum layer" of the minimal hierarchy means the derived spectrum, not
an independent input --- the referee-safe reading established by the spectrum-necessity audit
\cite{c05:QG295}.

\section{Summary}
\label{c05:sec:spectrum-summary}

This chapter has established the third member of the canonical architecture. The
actualization dynamics converge to a unique, content-independent attractor, the closure
fixed point (Theorem~\ref{c05:thm:attractor-unique}), whose size is the inevitable stable
fixed point $N=96$ (Theorem~\ref{c05:thm:n96-attractor}); the spectrum is the Laplacian
eigenspectrum of the converged network, its constants derived outputs
(Theorem~\ref{c05:thm:spectrum-eigenspectrum}). It is therefore unique, carries no choice,
and is not a primitive but the inevitable output required for physics
(Theorem~\ref{c05:thm:spectrum-not-primitive}). With the canonical core Difference →
Actualization → Spectrum complete (Theorem~\ref{c05:thm:core-completed}), physics is the
readout of the inevitable spectrum, which determines the observable space
(Theorem~\ref{c05:thm:no-physics-without-spectrum}); the physics chapters that follow derive
each sector as a read of this one spectrum.

\begin{thebibliography}{99}
\bibitem{c05:QG116} AT-QG Phase 116, \emph{Actualization Structures}, Docs/Research.
\bibitem{c05:QG159} AT-QG Phase 159, \emph{D96 Selection Origin}, Docs/Research.
\bibitem{c05:QG160} AT-QG Phase 160, \emph{Period-3 Seed Origin}, Docs/Research.
\bibitem{c05:QG272} AT-QG Phase 272, \emph{Sector Emergence Audit}, Docs/Research.
\bibitem{c05:QG282} AT-QG Phase 282, \emph{Boundary Origin Audit} (Closure Principle), Docs/Research.
\bibitem{c05:QG293} AT-QG Phase 293, \emph{Hierarchy Necessity Audit}, Docs/Research.
\bibitem{c05:QG294} AT-QG Phase 294, \emph{Minimal Theory Audit}, Docs/Research.
\bibitem{c05:QG295} AT-QG Phase 295, \emph{Spectrum Necessity Audit}, Docs/Research.
\bibitem{c05:QG296} AT-QG Phase 296, \emph{Reconstruction Audit}, Docs/Research.
\bibitem{c05:MONO007} AT-MONO007, \emph{Referee-Readiness Resolution}, Docs/Zenodo/Monograph.
\end{thebibliography}
 after the preamble
%         that defines the theorem environments (or include via the master file).
% ======================================================================

% ---- Theorem environments (define in the master preamble if not present) ----
% \newtheorem{definition}{Definition}[chapter]
% \newtheorem{lemma}[definition]{Lemma}
% \newtheorem{theorem}[definition]{Theorem}
% \newtheorem{corollary}[definition]{Corollary}
% \newtheorem{remark}[definition]{Remark}

\chapter{The Inevitable Spectrum}
\label{c05:chap:spectrum}

\section{Introduction}
\label{c05:sec:spectrum-introduction}

The canonical hierarchy of The Actualization Theory
\begin{equation}
\label{c05:eq:core}
\text{Difference}\;\longrightarrow\;\text{Actualization}\;\longrightarrow\;\text{Inevitable Spectrum}\;\longrightarrow\;\text{Physics}
\end{equation}
is complete as a derivation structure (Chapter~\ref{c04:chap:closure-minimal},
Theorem~\ref{c04:thm:minimal-hierarchy}). This chapter develops the third member of the
hierarchy: the \emph{inevitable spectrum}. We establish that the spectrum is not a primitive
of the theory but the \emph{derived output} of the actualization process: the dynamics
converge to a content-independent attractor of size $N=96$ within the accepted structural
class; the attractor is a stable fixed point, not a freely chosen input; and the spectrum is
the eigenspectrum of that attractor, carrying no independent choice
\cite{c05:QG116,c05:QG282,c05:QG295}.

We then establish the \emph{necessity} of the spectrum and the tested uniqueness of the
attractor within the canonical structural class: the attractor is content-independent
\cite{c05:QG116}, the size $N=96$ is inevitable within the accepted D96 selection
constraints --- not a choice \cite{c05:QG159,c05:QG160} --- and the spectrum is a
deterministic function of the converged network. This completes the canonical core
Difference → Actualization → Spectrum, whose first three members are now established as
primitive, derived, and inevitable output respectively; the physics chapters that follow
read all observable sectors from this one spectrum \cite{c05:QG293,c05:QG294,c05:QG296}.

\section{The Actualization Attractor}
\label{c05:sec:attractor}

\begin{theorem}[The actualization attractor: convergence, uniqueness, closure]
\label{c05:thm:attractor-unique}
The \emph{actualization attractor} is the converged state of the actualization process: the
configuration to which every initial pattern of Difference's count-producing dynamics
converges, with zero residual change. The actualization dynamics converge --- residual link
growth reaches zero and the topology saturates at a stable final state --- the attractor is
unique and content-independent (every initial activity pattern converges to the same final
geometry), and it is the closure fixed point of the theory.
\end{theorem}

\begin{proof}
The actualization-structures program \cite{c05:QG116} establishes that the actualization
dynamics converge with zero residual link growth: the self-reinforcing feedback (activity to
links to activity) is bounded by network capacity, so the topology saturates, and identical
link counts and identical span are reached from every initial pattern. Hence the attractor
is unique --- one final geometry, not a family --- and independent of the starting
configuration. By the Closure Principle (Chapter~\ref{c04:chap:closure-minimal},
Theorem~\ref{c04:thm:closure-holds}, Definition~\ref{c04:def:closure-principle}), the boundary
of the process is its stable fixed point; hence the unique attractor and the closure fixed
point coincide.
\end{proof}

\section{The Fixed Point \texorpdfstring{$N=96$}{N=96}}
\label{c05:sec:n96}

\begin{theorem}[The attractor is \texorpdfstring{$N=96$}{N=96}, not a choice; no freedom in the size]
\label{c05:thm:n96-attractor}
The \emph{$N=96$ attractor} is the converged network of the actualization process within the
canonical structural class: the stable fixed point whose size is $N=96$. The size is the
attractor of the dynamics, not a freely chosen input: the process converges to $N=96$ from
every initial pattern in the accepted comparison class, and no other tested size is a stable
fixed point of the dynamics.
\end{theorem}

\begin{proof}
The spectrum-necessity audit \cite{c05:QG295} establishes that $N=96$ is the attractor, not a
choice: the actualization flow has a stable fixed point at $N=96$, and the boundary is the
closure of that flow (the Closure Principle \cite{c05:QG282}). The D96-selection program adds
that $96$ is selected as the unique complete-Z2 natural size with the three-family octave
window; no other natural size satisfies these conditions \cite{c05:QG159,c05:QG160}. Hence
$N=96$ is the stable fixed point of the flow within the canonical comparison class,
content-independent (Theorem~\ref{c05:thm:attractor-unique}): an output of the dynamics, not
an input or free parameter.
\end{proof}

\subsection{Why \texorpdfstring{$N=96$}{N=96}?}
\label{c05:sec:why-n96}

\begin{table}[h]
\centering
\begin{tabular}{p{0.18\linewidth}p{0.28\linewidth}p{0.46\linewidth}}
\hline
\textbf{Argument class} & \textbf{What it uses} & \textbf{Status} \\
\hline
Selection & period-3 seed + Z2 half-shift; 3-family octave window; natural doubling chain $n=3\cdot2^k$; tested candidate set \cite{c05:QG159,c05:QG160} & D96 is the only rung in the tested structural class that satisfies all conditions \\
Necessity & closure of the actualization flow; stable fixed point; spectrum as eigenspectrum of the converged network \cite{c05:QG282,c05:QG295} & If the flow does not close, there is no spectrum layer and no physics readout \\
Uniqueness & content-independent attractor; same final geometry from every initial pattern \cite{c05:QG116} & Unique within the canonical comparison class examined \\
\hline
\end{tabular}
\end{table}

\paragraph{Why not 94, 95, 97?}
On the accepted selection criterion, these sizes fail the seed half-shift condition: the
period-3 seed requires $6 \mid n$ for the Z2 doublet structure, so $94,95,97$ are excluded
before the octave-window test is even applied. They are therefore not members of the
canonical candidate class used by QG159.

\paragraph{Why not 128?}
128 fails the same seed symmetry test ($128 \not\equiv 0 \pmod 6$) and, in the selection
audit, also fails the three-family window: its span is above the $[4,8)$ range, giving four
families rather than three \cite{c05:QG159}.

\paragraph{Exact status.}
The strongest referee-safe statement is: \emph{within the canonical structural constraints
actually used by AT --- period-3 seed, Z2 half-shift, three-family octave window, and the
tested comparison class --- $N=96$ is inevitable and uniquely selected.} A global proof that
no other closure could exist outside that structural class is \emph{not} claimed here.
That is the boundary.

\section{Spectrum Generation}
\label{c05:sec:spectrum-generation}

\begin{theorem}[The spectrum is the eigenspectrum of the attractor]
\label{c05:thm:spectrum-eigenspectrum}
\emph{Spectrum generation} is the derivation of the spectrum from the converged network: the
spectrum of the theory is the Laplacian eigenspectrum of the $N=96$ attractor --- $95$
positive modes with a single zero mode (the background) --- a deterministic function of the
attractor. Its spectral constants --- the moments $\Sigma m=95$,
$\Sigma\sqrt{m}=64.08$, $\Sigma m^2=229$, the occupancy moment $1900.25$, the span $6.40$,
and the doublet multiplicities $[42\times2,5,6]$ --- are outputs of the spectrum, hence of
the attractor.
\end{theorem}

\begin{proof}
The spectrum-necessity audit \cite{c05:QG295} establishes that the spectrum is the Laplacian
eigenspectrum of the converged network: the graph Laplacian of the $N=96$ attractor has
$95$ positive eigenvalues (the modes) and one zero eigenvalue (the background). Since the
attractor is unique and deterministic (Theorem~\ref{c05:thm:attractor-unique},
Theorem~\ref{c05:thm:n96-attractor}), the spectrum is a deterministic function of it --- an
output of the dynamics, not an additional input --- and its moments and structural features
are derived quantities, not primitive inputs. This is the canonical reading of the spectrum
layer of the minimal hierarchy \cite{c05:QG294}.
\end{proof}

\section{Uniqueness of the Spectrum}
\label{c05:sec:spectrum-uniqueness}

Uniqueness here means uniqueness \emph{within the canonical structural class}: the spectrum is
the eigenspectrum of the attractor (Theorem~\ref{c05:thm:attractor-unique},
Theorem~\ref{c05:thm:spectrum-eigenspectrum}), whose size is fixed at $N=96$
(Theorem~\ref{c05:thm:n96-attractor}), and the eigenspectrum of a given network is unique.
Hence there is exactly one spectrum of the theory in the accepted class --- one attractor,
one eigenspectrum --- and since the attractor carries no freedom, no spectral feature is
freely adjustable: the spectrum carries no choice at any level.

\section{Spectrum Necessity}
\label{c05:sec:spectrum-necessity}

\begin{theorem}[The spectrum is not primitive]
\label{c05:thm:spectrum-not-primitive}
The spectrum is not a primitive of the theory: it is the inevitable output of the
actualization attractor, carrying no choice and presupposing no independent input.
\end{theorem}

\begin{proof}
The spectrum-necessity audit \cite{c05:QG295} poses exactly this question --- is the spectrum
fundamental, or the inevitable fixed point of actualization? --- and establishes the latter.
The spectrum is the eigenspectrum of the converged network
(Theorem~\ref{c05:thm:spectrum-eigenspectrum}), which is the unique, fixed-size,
content-independent attractor (Theorem~\ref{c05:thm:attractor-unique},
Theorem~\ref{c05:thm:n96-attractor}); it therefore presupposes the actualization process and
its primitive Difference, and is an output, not an input. The reconstruction audit
\cite{c05:QG296} establishes that physics is the readout of the spectrum
(Chapter~\ref{c03:chap:actualization}, Corollary~\ref{c03:cor:physics-readout}), the sectors
being projection classes over the operator layer of the one spectrum.
\end{proof}

The spectrum is therefore also necessary: physics is the readout of the spectrum, so without
the spectrum there is no physics. The two properties are compatible: the spectrum is forced
by the dynamics (inevitable) and indispensable to the derivation (necessary).

\section{The Canonical Core Completed}
\label{c05:sec:core-completed}

\begin{theorem}[The canonical core is completed]
\label{c05:thm:core-completed}
The first three members of the canonical hierarchy are now established: Difference is the
primitive (Chapter~\ref{c01:chap:difference}), Actualization is the derived process face of
Difference (Chapter~\ref{c03:chap:actualization}), and the Spectrum is the inevitable output of
the actualization attractor (this chapter). The canonical core
\begin{equation}
\label{c05:eq:core-complete}
\text{Difference}\;\longrightarrow\;\text{Actualization}\;\longrightarrow\;\text{Inevitable Spectrum}
\end{equation}
is complete, anchoring the minimal hierarchy Difference → Actualization → Spectrum → Physics:
its first three members are established, and Physics is the spectrum readout developed in
the subsequent chapters.
\end{theorem}

\begin{proof}
The three members are established as primitive (Difference,
Theorem~\ref{c01:thm:minimal-primitives} of Chapter~\ref{c01:chap:difference}), derived process
(Actualization, Theorem~\ref{c03:thm:derived-not-primitive} of
Chapter~\ref{c03:chap:actualization}), and inevitable output (Spectrum,
Theorem~\ref{c05:thm:spectrum-not-primitive}), completing the first three members of the
minimal hierarchy \cite{c05:QG293,c05:QG294}. The fourth member, Physics, is the readout of the
spectrum (Chapter~\ref{c03:chap:actualization}, Corollary~\ref{c03:cor:physics-readout}),
developed in the physics chapters that follow. Every member of the hierarchy thus has a
canonical status: primitive, derived process, inevitable output, emergent readout.
\end{proof}

\section{Consequences for Physics}
\label{c05:sec:physics-consequences}

All observable physics is the readout of the inevitable spectrum: the resonance structure
projected by the operator layer, from which the sectors of physics emerge. The spectrum is
the unique, inevitable output of the attractor (Theorem~\ref{c05:thm:spectrum-not-primitive});
the sectors of physics are projection classes over the operator layer of the one spectrum
(Chapter~\ref{c03:chap:actualization}, Theorem~\ref{c03:thm:resonance-readout}). Hence every
sector --- masses, couplings, mixings, gravity, cosmology --- is a projection of the one
actualization-driven spectrum \cite{c05:QG272}.

\begin{theorem}[No physics without the spectrum]
\label{c05:thm:no-physics-without-spectrum}
There is no physics in the theory without the spectrum: every derived observable is a
function of the spectral structure, and no observable bypasses the spectrum layer. Hence the
observable space of the theory is determined by the inevitable spectrum: the set of derived
observables is the set of reads of the unique spectrum.
\end{theorem}

\begin{proof}
The reconstruction audit \cite{c05:QG296} reconstructs the major results of the theory from the
minimal hierarchy, with the spectrum layer as the source of the spectral constants used by
all sectors; hence every observable passes through the spectrum layer, and there is no
physics without it. Since the attractor is unique (Theorem~\ref{c05:thm:attractor-unique}),
the spectrum is unique; every observable is a read of the one inevitable spectrum, and the
observable space is the set of those reads.
\end{proof}

Consistent with MONO007 \cite{c05:MONO007}, the spectrum is never presented as a primitive in
this chapter: every statement positions it as the derived output of the actualization
attractor, and the "spectrum layer" of the minimal hierarchy means the derived spectrum, not
an independent input --- the referee-safe reading established by the spectrum-necessity audit
\cite{c05:QG295}.

\section{Summary}
\label{c05:sec:spectrum-summary}

This chapter has established the third member of the canonical architecture. The
actualization dynamics converge to a unique, content-independent attractor, the closure
fixed point (Theorem~\ref{c05:thm:attractor-unique}), whose size is the inevitable stable
fixed point $N=96$ (Theorem~\ref{c05:thm:n96-attractor}); the spectrum is the Laplacian
eigenspectrum of the converged network, its constants derived outputs
(Theorem~\ref{c05:thm:spectrum-eigenspectrum}). It is therefore unique, carries no choice,
and is not a primitive but the inevitable output required for physics
(Theorem~\ref{c05:thm:spectrum-not-primitive}). With the canonical core Difference →
Actualization → Spectrum complete (Theorem~\ref{c05:thm:core-completed}), physics is the
readout of the inevitable spectrum, which determines the observable space
(Theorem~\ref{c05:thm:no-physics-without-spectrum}); the physics chapters that follow derive
each sector as a read of this one spectrum.

\begin{thebibliography}{99}
\bibitem{c05:QG116} AT-QG Phase 116, \emph{Actualization Structures}, Docs/Research.
\bibitem{c05:QG159} AT-QG Phase 159, \emph{D96 Selection Origin}, Docs/Research.
\bibitem{c05:QG160} AT-QG Phase 160, \emph{Period-3 Seed Origin}, Docs/Research.
\bibitem{c05:QG272} AT-QG Phase 272, \emph{Sector Emergence Audit}, Docs/Research.
\bibitem{c05:QG282} AT-QG Phase 282, \emph{Boundary Origin Audit} (Closure Principle), Docs/Research.
\bibitem{c05:QG293} AT-QG Phase 293, \emph{Hierarchy Necessity Audit}, Docs/Research.
\bibitem{c05:QG294} AT-QG Phase 294, \emph{Minimal Theory Audit}, Docs/Research.
\bibitem{c05:QG295} AT-QG Phase 295, \emph{Spectrum Necessity Audit}, Docs/Research.
\bibitem{c05:QG296} AT-QG Phase 296, \emph{Reconstruction Audit}, Docs/Research.
\bibitem{c05:MONO007} AT-MONO007, \emph{Referee-Readiness Resolution}, Docs/Zenodo/Monograph.
\end{thebibliography}
 after the preamble
%         that defines the theorem environments (or include via the master file).
% ======================================================================

% ---- Theorem environments (define in the master preamble if not present) ----
% \newtheorem{definition}{Definition}[chapter]
% \newtheorem{lemma}[definition]{Lemma}
% \newtheorem{theorem}[definition]{Theorem}
% \newtheorem{corollary}[definition]{Corollary}
% \newtheorem{remark}[definition]{Remark}

\chapter{The Inevitable Spectrum}
\label{c05:chap:spectrum}

\section{Introduction}
\label{c05:sec:spectrum-introduction}

The canonical hierarchy of The Actualization Theory
\begin{equation}
\label{c05:eq:core}
\text{Difference}\;\longrightarrow\;\text{Actualization}\;\longrightarrow\;\text{Inevitable Spectrum}\;\longrightarrow\;\text{Physics}
\end{equation}
is complete as a derivation structure (Chapter~\ref{c04:chap:closure-minimal},
Theorem~\ref{c04:thm:minimal-hierarchy}). This chapter develops the third member of the
hierarchy: the \emph{inevitable spectrum}. We establish that the spectrum is not a primitive
of the theory but the \emph{derived output} of the actualization process: the dynamics
converge to a content-independent attractor of size $N=96$ within the accepted structural
class; the attractor is a stable fixed point, not a freely chosen input; and the spectrum is
the eigenspectrum of that attractor, carrying no independent choice
\cite{c05:QG116,c05:QG282,c05:QG295}.

We then establish the \emph{necessity} of the spectrum and the tested uniqueness of the
attractor within the canonical structural class: the attractor is content-independent
\cite{c05:QG116}, the size $N=96$ is inevitable within the accepted D96 selection
constraints --- not a choice \cite{c05:QG159,c05:QG160} --- and the spectrum is a
deterministic function of the converged network. This completes the canonical core
Difference → Actualization → Spectrum, whose first three members are now established as
primitive, derived, and inevitable output respectively; the physics chapters that follow
read all observable sectors from this one spectrum \cite{c05:QG293,c05:QG294,c05:QG296}.

\section{The Actualization Attractor}
\label{c05:sec:attractor}

\begin{theorem}[The actualization attractor: convergence, uniqueness, closure]
\label{c05:thm:attractor-unique}
The \emph{actualization attractor} is the converged state of the actualization process: the
configuration to which every initial pattern of Difference's count-producing dynamics
converges, with zero residual change. The actualization dynamics converge --- residual link
growth reaches zero and the topology saturates at a stable final state --- the attractor is
unique and content-independent (every initial activity pattern converges to the same final
geometry), and it is the closure fixed point of the theory.
\end{theorem}

\begin{proof}
The actualization-structures program \cite{c05:QG116} establishes that the actualization
dynamics converge with zero residual link growth: the self-reinforcing feedback (activity to
links to activity) is bounded by network capacity, so the topology saturates, and identical
link counts and identical span are reached from every initial pattern. Hence the attractor
is unique --- one final geometry, not a family --- and independent of the starting
configuration. By the Closure Principle (Chapter~\ref{c04:chap:closure-minimal},
Theorem~\ref{c04:thm:closure-holds}, Definition~\ref{c04:def:closure-principle}), the boundary
of the process is its stable fixed point; hence the unique attractor and the closure fixed
point coincide.
\end{proof}

\section{The Fixed Point \texorpdfstring{$N=96$}{N=96}}
\label{c05:sec:n96}

\begin{theorem}[The attractor is \texorpdfstring{$N=96$}{N=96}, not a choice; no freedom in the size]
\label{c05:thm:n96-attractor}
The \emph{$N=96$ attractor} is the converged network of the actualization process within the
canonical structural class: the stable fixed point whose size is $N=96$. The size is the
attractor of the dynamics, not a freely chosen input: the process converges to $N=96$ from
every initial pattern in the accepted comparison class, and no other tested size is a stable
fixed point of the dynamics.
\end{theorem}

\begin{proof}
The spectrum-necessity audit \cite{c05:QG295} establishes that $N=96$ is the attractor, not a
choice: the actualization flow has a stable fixed point at $N=96$, and the boundary is the
closure of that flow (the Closure Principle \cite{c05:QG282}). The D96-selection program adds
that $96$ is selected as the unique complete-Z2 natural size with the three-family octave
window; no other natural size satisfies these conditions \cite{c05:QG159,c05:QG160}. Hence
$N=96$ is the stable fixed point of the flow within the canonical comparison class,
content-independent (Theorem~\ref{c05:thm:attractor-unique}): an output of the dynamics, not
an input or free parameter.
\end{proof}

\subsection{Why \texorpdfstring{$N=96$}{N=96}?}
\label{c05:sec:why-n96}

\begin{table}[h]
\centering
\begin{tabular}{p{0.18\linewidth}p{0.28\linewidth}p{0.46\linewidth}}
\hline
\textbf{Argument class} & \textbf{What it uses} & \textbf{Status} \\
\hline
Selection & period-3 seed + Z2 half-shift; 3-family octave window; natural doubling chain $n=3\cdot2^k$; tested candidate set \cite{c05:QG159,c05:QG160} & D96 is the only rung in the tested structural class that satisfies all conditions \\
Necessity & closure of the actualization flow; stable fixed point; spectrum as eigenspectrum of the converged network \cite{c05:QG282,c05:QG295} & If the flow does not close, there is no spectrum layer and no physics readout \\
Uniqueness & content-independent attractor; same final geometry from every initial pattern \cite{c05:QG116} & Unique within the canonical comparison class examined \\
\hline
\end{tabular}
\end{table}

\paragraph{Why not 94, 95, 97?}
On the accepted selection criterion, these sizes fail the seed half-shift condition: the
period-3 seed requires $6 \mid n$ for the Z2 doublet structure, so $94,95,97$ are excluded
before the octave-window test is even applied. They are therefore not members of the
canonical candidate class used by QG159.

\paragraph{Why not 128?}
128 fails the same seed symmetry test ($128 \not\equiv 0 \pmod 6$) and, in the selection
audit, also fails the three-family window: its span is above the $[4,8)$ range, giving four
families rather than three \cite{c05:QG159}.

\paragraph{Exact status.}
The strongest referee-safe statement is: \emph{within the canonical structural constraints
actually used by AT --- period-3 seed, Z2 half-shift, three-family octave window, and the
tested comparison class --- $N=96$ is inevitable and uniquely selected.} A global proof that
no other closure could exist outside that structural class is \emph{not} claimed here.
That is the boundary.

\section{Spectrum Generation}
\label{c05:sec:spectrum-generation}

\begin{theorem}[The spectrum is the eigenspectrum of the attractor]
\label{c05:thm:spectrum-eigenspectrum}
\emph{Spectrum generation} is the derivation of the spectrum from the converged network: the
spectrum of the theory is the Laplacian eigenspectrum of the $N=96$ attractor --- $95$
positive modes with a single zero mode (the background) --- a deterministic function of the
attractor. Its spectral constants --- the moments $\Sigma m=95$,
$\Sigma\sqrt{m}=64.08$, $\Sigma m^2=229$, the occupancy moment $1900.25$, the span $6.40$,
and the doublet multiplicities $[42\times2,5,6]$ --- are outputs of the spectrum, hence of
the attractor.
\end{theorem}

\begin{proof}
The spectrum-necessity audit \cite{c05:QG295} establishes that the spectrum is the Laplacian
eigenspectrum of the converged network: the graph Laplacian of the $N=96$ attractor has
$95$ positive eigenvalues (the modes) and one zero eigenvalue (the background). Since the
attractor is unique and deterministic (Theorem~\ref{c05:thm:attractor-unique},
Theorem~\ref{c05:thm:n96-attractor}), the spectrum is a deterministic function of it --- an
output of the dynamics, not an additional input --- and its moments and structural features
are derived quantities, not primitive inputs. This is the canonical reading of the spectrum
layer of the minimal hierarchy \cite{c05:QG294}.
\end{proof}

\section{Uniqueness of the Spectrum}
\label{c05:sec:spectrum-uniqueness}

Uniqueness here means uniqueness \emph{within the canonical structural class}: the spectrum is
the eigenspectrum of the attractor (Theorem~\ref{c05:thm:attractor-unique},
Theorem~\ref{c05:thm:spectrum-eigenspectrum}), whose size is fixed at $N=96$
(Theorem~\ref{c05:thm:n96-attractor}), and the eigenspectrum of a given network is unique.
Hence there is exactly one spectrum of the theory in the accepted class --- one attractor,
one eigenspectrum --- and since the attractor carries no freedom, no spectral feature is
freely adjustable: the spectrum carries no choice at any level.

\section{Spectrum Necessity}
\label{c05:sec:spectrum-necessity}

\begin{theorem}[The spectrum is not primitive]
\label{c05:thm:spectrum-not-primitive}
The spectrum is not a primitive of the theory: it is the inevitable output of the
actualization attractor, carrying no choice and presupposing no independent input.
\end{theorem}

\begin{proof}
The spectrum-necessity audit \cite{c05:QG295} poses exactly this question --- is the spectrum
fundamental, or the inevitable fixed point of actualization? --- and establishes the latter.
The spectrum is the eigenspectrum of the converged network
(Theorem~\ref{c05:thm:spectrum-eigenspectrum}), which is the unique, fixed-size,
content-independent attractor (Theorem~\ref{c05:thm:attractor-unique},
Theorem~\ref{c05:thm:n96-attractor}); it therefore presupposes the actualization process and
its primitive Difference, and is an output, not an input. The reconstruction audit
\cite{c05:QG296} establishes that physics is the readout of the spectrum
(Chapter~\ref{c03:chap:actualization}, Corollary~\ref{c03:cor:physics-readout}), the sectors
being projection classes over the operator layer of the one spectrum.
\end{proof}

The spectrum is therefore also necessary: physics is the readout of the spectrum, so without
the spectrum there is no physics. The two properties are compatible: the spectrum is forced
by the dynamics (inevitable) and indispensable to the derivation (necessary).

\section{The Canonical Core Completed}
\label{c05:sec:core-completed}

\begin{theorem}[The canonical core is completed]
\label{c05:thm:core-completed}
The first three members of the canonical hierarchy are now established: Difference is the
primitive (Chapter~\ref{c01:chap:difference}), Actualization is the derived process face of
Difference (Chapter~\ref{c03:chap:actualization}), and the Spectrum is the inevitable output of
the actualization attractor (this chapter). The canonical core
\begin{equation}
\label{c05:eq:core-complete}
\text{Difference}\;\longrightarrow\;\text{Actualization}\;\longrightarrow\;\text{Inevitable Spectrum}
\end{equation}
is complete, anchoring the minimal hierarchy Difference → Actualization → Spectrum → Physics:
its first three members are established, and Physics is the spectrum readout developed in
the subsequent chapters.
\end{theorem}

\begin{proof}
The three members are established as primitive (Difference,
Theorem~\ref{c01:thm:minimal-primitives} of Chapter~\ref{c01:chap:difference}), derived process
(Actualization, Theorem~\ref{c03:thm:derived-not-primitive} of
Chapter~\ref{c03:chap:actualization}), and inevitable output (Spectrum,
Theorem~\ref{c05:thm:spectrum-not-primitive}), completing the first three members of the
minimal hierarchy \cite{c05:QG293,c05:QG294}. The fourth member, Physics, is the readout of the
spectrum (Chapter~\ref{c03:chap:actualization}, Corollary~\ref{c03:cor:physics-readout}),
developed in the physics chapters that follow. Every member of the hierarchy thus has a
canonical status: primitive, derived process, inevitable output, emergent readout.
\end{proof}

\section{Consequences for Physics}
\label{c05:sec:physics-consequences}

All observable physics is the readout of the inevitable spectrum: the resonance structure
projected by the operator layer, from which the sectors of physics emerge. The spectrum is
the unique, inevitable output of the attractor (Theorem~\ref{c05:thm:spectrum-not-primitive});
the sectors of physics are projection classes over the operator layer of the one spectrum
(Chapter~\ref{c03:chap:actualization}, Theorem~\ref{c03:thm:resonance-readout}). Hence every
sector --- masses, couplings, mixings, gravity, cosmology --- is a projection of the one
actualization-driven spectrum \cite{c05:QG272}.

\begin{theorem}[No physics without the spectrum]
\label{c05:thm:no-physics-without-spectrum}
There is no physics in the theory without the spectrum: every derived observable is a
function of the spectral structure, and no observable bypasses the spectrum layer. Hence the
observable space of the theory is determined by the inevitable spectrum: the set of derived
observables is the set of reads of the unique spectrum.
\end{theorem}

\begin{proof}
The reconstruction audit \cite{c05:QG296} reconstructs the major results of the theory from the
minimal hierarchy, with the spectrum layer as the source of the spectral constants used by
all sectors; hence every observable passes through the spectrum layer, and there is no
physics without it. Since the attractor is unique (Theorem~\ref{c05:thm:attractor-unique}),
the spectrum is unique; every observable is a read of the one inevitable spectrum, and the
observable space is the set of those reads.
\end{proof}

Consistent with MONO007 \cite{c05:MONO007}, the spectrum is never presented as a primitive in
this chapter: every statement positions it as the derived output of the actualization
attractor, and the "spectrum layer" of the minimal hierarchy means the derived spectrum, not
an independent input --- the referee-safe reading established by the spectrum-necessity audit
\cite{c05:QG295}.

\section{Summary}
\label{c05:sec:spectrum-summary}

This chapter has established the third member of the canonical architecture. The
actualization dynamics converge to a unique, content-independent attractor, the closure
fixed point (Theorem~\ref{c05:thm:attractor-unique}), whose size is the inevitable stable
fixed point $N=96$ (Theorem~\ref{c05:thm:n96-attractor}); the spectrum is the Laplacian
eigenspectrum of the converged network, its constants derived outputs
(Theorem~\ref{c05:thm:spectrum-eigenspectrum}). It is therefore unique, carries no choice,
and is not a primitive but the inevitable output required for physics
(Theorem~\ref{c05:thm:spectrum-not-primitive}). With the canonical core Difference →
Actualization → Spectrum complete (Theorem~\ref{c05:thm:core-completed}), physics is the
readout of the inevitable spectrum, which determines the observable space
(Theorem~\ref{c05:thm:no-physics-without-spectrum}); the physics chapters that follow derive
each sector as a read of this one spectrum.

\begin{thebibliography}{99}
\bibitem{c05:QG116} AT-QG Phase 116, \emph{Actualization Structures}, Docs/Research.
\bibitem{c05:QG159} AT-QG Phase 159, \emph{D96 Selection Origin}, Docs/Research.
\bibitem{c05:QG160} AT-QG Phase 160, \emph{Period-3 Seed Origin}, Docs/Research.
\bibitem{c05:QG272} AT-QG Phase 272, \emph{Sector Emergence Audit}, Docs/Research.
\bibitem{c05:QG282} AT-QG Phase 282, \emph{Boundary Origin Audit} (Closure Principle), Docs/Research.
\bibitem{c05:QG293} AT-QG Phase 293, \emph{Hierarchy Necessity Audit}, Docs/Research.
\bibitem{c05:QG294} AT-QG Phase 294, \emph{Minimal Theory Audit}, Docs/Research.
\bibitem{c05:QG295} AT-QG Phase 295, \emph{Spectrum Necessity Audit}, Docs/Research.
\bibitem{c05:QG296} AT-QG Phase 296, \emph{Reconstruction Audit}, Docs/Research.
\bibitem{c05:MONO007} AT-MONO007, \emph{Referee-Readiness Resolution}, Docs/Zenodo/Monograph.
\end{thebibliography}
